\exercisegroup

In all the Exercises for § 5, \(\mathfrak{C}\) denotes an equalitarian theory.

\begin{enumerate}
\def\labelenumi{\arabic{enumi}.}
\item
  The relation \(x = y\) is functional in \(x\) in \(\mathfrak{C}\) .
\item
  Let R be a relation in G and let x, y be distinct letters. Then the relations \((\exists x)(x = y \text{ and } R)\) , \((y|x)R\) are equivalent in G.
\item
  Let R and S be relations in C, let T be a term in C, and let x and y be distinct letters. Suppose that y is not a constant of T and that x does not appear in T. Let \(C'\) be the theory obtained by adjoining S to the axioms of C. If R is functional in x in \(C'\) , and if \((T|y)S\) is a theorem in C, then the relation \((T|y)R\) is functional in X in C.
\item
  Let R and S be relations in C, and let x be a letter which is not a constant of C. If R is functional in x in C, and if \(R \leftrightarrow S\) is a theorem in C, then S is functional in x in C.
\item
  Let \(\pmb{R}\) , \(\pmb{S}\) , \(\pmb{T}\) be relations in \(\mathfrak{C}\) and let \(\pmb{x}\) be a letter. If \(\pmb{R}\) is functional in \(\pmb{x}\) in \(\mathfrak{C}\) , show that the following relations are theorems in \(\mathfrak{C}\) :
\end{enumerate}

\[
(\text { not   } (\exists x) (R \text {   and   } S)) \iff (\exists x) (R \text {   and   } (\text { not   } S)),
\]

\((\exists x)(R \text{ and } (S \text{ and } T)) \Longleftrightarrow (\exists x)(R \text{ and } S) \text{ and } (\exists x)(R \text{ and } T)),\) \((\exists x)(R \text{ and } (S \text{ or } T)) \Longleftrightarrow ((\exists x)(R \text{ and } S) \text{ or } (\exists x)(R \text{ and } T)).\)

\begin{enumerate}
\def\labelenumi{\arabic{enumi}.}
\setcounter{enumi}{5}
\item
  Show that if the scheme \((\exists x)R \Longrightarrow R\) provides implicit axioms in \(\mathfrak{C}\) , then \(x = y\) is a theorem in \(\mathfrak{C}\) (cf.~Exercise 1).
\item
  Show that if the scheme \((\pmb{R} \Longleftrightarrow \pmb{S}) \Longrightarrow (\tau_{\pmb{x}}(\pmb{R}) = \tau_{\pmb{x}}(\pmb{S}))\) provides implicit axioms in \(\mathfrak{C}\) , then \(x = y\) is a theorem in \(\mathfrak{C}\) . (Take \(\pmb{R}\) to be the relation \(x = x\) , \(\pmb{S}\) the relation \(x = y\) , and then substitute \(x\) for \(y\) in the axiom thus obtained.)\footnote{We shall see that this theory admits no model in the theory of sets (§ 4, no. 8) having at least two elements (i.e. in which the relation \((\exists x)(\exists y)(x \neq y)\) is true).}
\end{enumerate}

